% --- 标题 ---
\title{\timu{报告骨架}}
\author{测试}
\date{\today}

\maketitle
\thispagestyle{empty}
\newpage

% --- 目录（需编译两次才能生成） ---
\tableofcontents
\newpage

% ============================================================
% 可用字体命令:
%   \timu{...}         — 一级标题字体
%   \yijitimu{...}     — 二级标题字体
%   \erjitimu{...}     — 三级标题字体
%   \zz{...}           — 正文强调
%   \dw{...}           — 正文引用
%   \seeref{...}       — 参见引用（带 ▶ 符号）
%
% 可用环境:
%   \begin{mylist}...\end{mylist}       — 多级智能列表
%   \begin{mycolumns}[n]...\end{mycolumns} — n 栏分栏
%
% 表格与图片:
%   \begin{tabularx}{\linewidth}{|c|X|}... — 自适应表格
%   \begin{figure}[H] \includegraphics[...]{...} — 图片插入
% ============================================================

\section{项目概述}

\subsection{背景}

% 说明项目背景、目标和范围。

\subsection{核心需求}

% 在此处列出核心需求：
%
% \begin{mylist}
%   \item 需求一
%   \item 需求二
%   \item 需求三
% \end{mylist}

\section{技术方案}

\subsection{架构设计}

% 描述系统架构。

% 下表展示技术选型：

\begin{table}[H]
\centering
\caption{技术选型}
\begin{tabularx}{\linewidth}{|l|X|}
\hline
\textbf{组件} & \textbf{选型} \\ \hline
后端框架 & Spring Boot / FastAPI \\ \hline
数据库 & PostgreSQL \\ \hline
消息队列 & Kafka \\ \hline
\end{tabularx}
\end{table}

\subsection{实施计划}

% 可按阶段划分：
%
% \begin{mylist}
%   \item 阶段一（第1-4周）：基础设施搭建
%   \item 阶段二（第5-12周）：核心功能开发
%   \item 阶段三（第13-16周）：测试部署
% \end{mylist}

\section{风险评估}

\begin{table}[H]
\centering
\caption{风险登记表}
\begin{tabularx}{\linewidth}{|c|X|c|X|}
\hline
\textbf{等级} & \textbf{风险描述} & \textbf{概率} & \textbf{应对措施} \\ \hline
高 & 技术栈学习成本 & 中 & 提前培训 \\ \hline
中 & 需求变更 & 高 & 敏捷迭代 \\ \hline
低 & 人员变动 & 低 & 文档沉淀 \\ \hline
\end{tabularx}
\end{table}

\section{总结}

% 总结报告的结论和建议。

